%% For double-blind review submission, w/o author-identifying information,
%% add 'anonymous' to the document class.
%%
%% For single-blind review submission, add the [review] option to
%% the document class
%%
%% For final camera ready submission, add 'final' to the document class.
%%
\documentclass[sigplan,final,acmsmall,review,screen]{acmart}

%% Conference information
\settopmatter{printfolios=true}

%% Bibliography style
\bibliographystyle{ACM-Reference-Format}
%% Citation style
\citestyle{acmauthoryear}

\usepackage{listings}
\usepackage{xcolor}
\usepackage{booktabs}
\usepackage{placeins}
\usepackage{float}
\usepackage{tcolorbox}

% Code listings styling
\lstset{
  basicstyle=\ttfamily\small,
  keywordstyle=\color{blue},
  commentstyle=\color{gray},
  stringstyle=\color{red},
  breaklines=true,
  postbreak=\mbox{\textcolor{red}{$\hookrightarrow$}\space},
  numbers=left,
  numberstyle=\tiny,
  stepnumber=5,
  columns=fullflexible
}

% Theorem box styling
\newtcolorbox{boxedtheorem}{colback=blue!5!white,colframe=blue!75!black,title=Theorem,fonttitle=\bfseries,arc=0mm}

\begin{document}

%% Title information
\title{Object-Level Dispatch Caching is Counterproductive: \\Why 99.99\% Hit Rates Fail to Improve Performance}

%%\author{Frank Adrian}
%%\affiliation{%
%%  \institution{Ancar Technology}
%%%%  \country{United States}
%%}
%%\email{frank.adrian314159@gmail.com}

\begin{abstract}
\textbf{Abstract.}~We establish an analytic cost model demonstrating that object-level dispatch caching is fundamentally limited by an irreducible efficiency gap:
  cache lookup cost (20--300~ns) cannot simultaneously beat both ultra-fast dispatch (1--100~ns) and expensive dispatch (>10~µs). This gap is rooted in CPU physics: memory access latency and hash computation form an irreducible lower bound.

  We validate this across twenty-four language implementations (compiled, JIT, interpreted), showing object-level caching fails despite 99.99\%+ hit rates. Failure is universal across 6 language families and 5 dispatch mechanisms: modern implementations optimize dispatch below cache lookup cost, making caching counterproductive. Results: 21/24 implementations fail (>1.1× slowdown).

  The paper makes three core contributions: (1)~an analytic bound characterizing the irreducible gap between cache lookup and dispatch costs, explaining why
  this optimization fails across all contemporary implementations; (2)~empirical validation across twenty-four diverse implementations, showing the gap is
  fundamental rather than language-specific; (3)~design space analysis identifying theoretical conditions where caching becomes viable (expensive predicates,
  lazy JIT startup, polyglot dispatch, deliberate semantic richness in dispatch), explaining why no current language exhibits such conditions by design.
  A critical distinction emerges: caching effectiveness (hit rate) is orthogonal to caching efficiency (overhead per call), explaining why high hit rates
  do not predict performance improvements.
\end{abstract}

\begin{CCSXML}
<ccs2012>
<concept>
<concept_id>10011007.10011006.10011050.10011017</concept_id>
<concept_desc>Software and its engineering~Just-in-time compilers</concept_desc>
<concept_keyword>inline caching</concept_keyword>
</concept>
<concept>
<concept_id>10011007.10011006.10011050.10011066</concept_id>
<concept_desc>Software and its engineering~Dynamic compilers</concept_desc>
<concept_keyword>dispatch optimization</concept_keyword>
</concept>
<concept>
<concept_id>10011007.10011006.10011041.10011047</concept_id>
<concept_desc>Software and its engineering~Source code generation</concept_desc>
<concept_keyword>performance</concept_keyword>
</concept>
</ccs2012>
\end{CCSXML}

\ccsdesc[500]{Software and its engineering~Just-in-time compilers}
\ccsdesc[300]{Software and its engineering~Dynamic compilers}
\ccsdesc[300]{Software and its engineering~Source code generation}

\keywords{inline caching, dispatch optimization, performance analysis, dynamic languages, universality}

\maketitle

\section{Introduction}

Efficient method dispatch is critical for object-oriented and dynamically-typed languages.
Polymorphic inline caching (PIC)~\cite{hoelzle1991optimizing} has emerged as the standard
optimization approach across modern dynamic language implementations. JavaScript engines like V8,
Python's PyPy, and the GraalVM achieve substantial speedups through inline caches that store
recently-seen type combinations, avoiding expensive dispatch predicates on subsequent calls.

The success of inline caching in V8, PyPy, and GraalVM has motivated decades of research into caching strategies. However, these
successes rely on \emph{machine-code inline caching}: JIT compilers embed type checks and direct jumps into compiled code, achieving near-zero overhead through specialization. Why, then, evaluate an approach that JITs have already solved? Because developers and language designers still mistakenly reach for object-level caching—a real intuition trap with practical consequences.

\textbf{Real-World Examples}: Python developers use \texttt{@lru\_cache} on dispatchers (overhead exceeds type-check cost). Clojure multimethods include caching but show no speedup on repeated types. These production patterns show developers reasoning correctly about caching in isolation but failing to account for modern optimization. A critical question: \textbf{is there a fundamental gap between cache lookup cost and dispatch cost savings}? Can cache lookup latency simultaneously beat ultra-fast dispatch and the expensive dispatch (>10~µs) theoretically required for break-even?

  To answer this question, we develop an analytic cost model characterizing the irreducible efficiency gap. Cache lookup cost is bounded below
  by memory access latency and hash computation (20--300~ns), while modern language implementations optimize dispatch to costs below this threshold
  (1--100~ns). This gap is not language-specific but a consequence of CPU physics. We validate this model through comprehensive empirical study
  across twenty-four language implementations spanning six language families and five dispatch mechanisms, showing that multi-threaded lock contention turns this failure into a catastrophic one (up to 88$\times$ slowdown).

  \textbf{Core finding}: The analytic cost model predicts universal failure across all contemporary implementations optimizing dispatch for speed.
Empirical validation confirms this: 21/24 implementations show clear failure (>1.1× slowdown) despite achieving 99.99\%+ cache hit rates.
The only implementation approaching break-even (CCL) does so when its baseline dispatch cost matches cache lookup latency,
validating the theory. Design space analysis identifies conditions where caching becomes viable (expensive predicates, lazy JIT, polyglot dispatch),
but no current language implements such designs, confirming convergence on speed-optimized dispatch.

\subsection{Core Finding}

Testing across twenty-four implementations and five dispatch mechanisms reveals:
\begin{itemize}
  \item \textbf{21/24 implementations}: Clear failure (>1.1× slowdown)
  \item \textbf{2/24 implementations}: Marginal/near break-even (1.02× CCL, 0.99× Racket; within noise)
  \item \textbf{1/24 implementations}: Actual speedup (0.77× Python match; genuine improvement)
  \item \textbf{4/5 mechanisms}: Clear failure (slowdown >1.1×) across all implementations
  \item \textbf{1/5 mechanisms}: At break-even (hash dispatch, neutral performance; 5\% marginal within noise—no speedup, not a win)
\end{itemize}

Failure severity spans three orders of magnitude:
\begin{itemize}
  \item Marginal (1.15× slowdown in Go, 1.10× in Typed Racket)
  \item Severe (5.3× slowdown in SBCL, 84× in LuaJIT)
  \item Catastrophic (V8, C2 with $\infty \times$ slowdown from escape analysis defeating object allocation)
\end{itemize}

\textbf{Interpretation}: The two implementations approaching break-even (CCL, Racket) share a common property: relatively expensive
baseline dispatch (360~ns for CCL, 420~ns for Racket). This suggests that languages with naturally slower dispatch might benefit from caching,
a pattern we explore in depth in Section~\ref{sec:patterns}.

\subsection{Contributions}

\begin{enumerate}
  \item \textbf{Analytic cost model proving irreducible gap}. We establish a performance bound showing that caching fails when
    cache lookup cost $C$ (20--300~ns, determined by memory latency and hash computation) cannot be simultaneously less than both
    (a)~ultra-fast optimized dispatch $F_{\text{opt}}$ (1--100~ns, achievable through compiler specialization) and (b)~expensive dispatch
    $F_{\text{expensive}}$ (>10~µs, required for break-even but non-existent in practice). This gap is fundamental, rooted in CPU physics
    (memory hierarchy and arithmetic latency) rather than implementation choices. Observation: No language optimizing dispatch can simultaneously
    achieve both low baseline cost (<100~ns) and benefit from object-level caching.

  \item \textbf{Comprehensive empirical validation across 24 diverse implementations}. We test compiled natives (SBCL, CCL, LispWorks, Chez, Go),
    JVM-based languages (ABCL, Clojure), bytecode/compiled/tracing JITs (C2, GraalVM, Dart, V8, Julia, LuaJIT, PyPy), optional types (Typed Racket, TypeScript), and interpreters
    (Python, Ruby, Lua 5.1, Racket)—spanning six language families and three execution models. Results confirm the irreducible gap: 21/24 implementations show clear failure
    (>1.1× slowdown) despite 99.99\%+ hit rates, demonstrating the gap is systematic (rooted in universal CPU physics) rather than language-specific.

  \item \textbf{Critical distinction: effectiveness vs. efficiency}. High cache hit rates (99.99\%+) do not predict performance improvement.
    Caching effectiveness (hit rate) is orthogonal to caching efficiency (per-call overhead). This explains the apparent paradox: universal cache success
    in theory (near-perfect hit rates) but universal failure in practice (overhead exceeds savings).

  \item \textbf{Design space characterization}. We identify theoretical conditions where caching becomes viable: (1)~expensive predicates (>1~µs dispatch cost; empirically validated via regex: 9--35× speedup),
    (2)~lazy JIT compilation (cold-start overhead; GraalVM: 1.39× speedup), (3)~polyglot dispatch (FFI boundaries; real C FFI: 1.74×, JNI: 1.68×),
    (4)~deliberate semantic richness in dispatch (e.g., Clojure multimethods with expensive predicates). No current language implements such designs by default, explaining universal failure. This reveals a convergence property: modern language implementations
    have independently optimized dispatch to the speed-optimal regime where caching fails.

  \item \textbf{Pattern analysis by dispatch cost}. Caching viability increases monotonically with baseline dispatch cost until overhead fraction becomes
    negligible. Break-even occurs when cache lookup cost (~8--30~ns) matches baseline dispatch cost, validated by hash dispatch at 5% margin and
    CCL/Racket at near break-even (1.02×/0.99×). This pattern reveals why failure is universal: modern implementations optimize to the regime
    where cache lookup cost dominates.
\end{enumerate}

\subsection{Theorem 1: Dispatch Caching Failure in Speed-Optimized Systems}

\textbf{Statement}: Speed-optimized JIT systems (V8, C2) defeat caching via specialization. Yet conservative compiled systems (SBCL) preserve expensive predicates. Example: V8 optimizes constraint solver 5323~ns → 6.69~ns (0.00×), SBCL keeps it 22,833~ns → 674× speedup. Caching fails universally only for speed-optimized JIT implementations; design space remains viable for conservative compiled languages.

\section{Background}

\subsection{Polymorphic Inline Caching in Dynamic Languages}

Inline caching was introduced by Hölzle, Chambers, and Ungar~\cite{hoelzle1991optimizing}
for Smalltalk to address the overhead of method dispatch on every call. The core idea is to
store (type-tuple, target-method) pairs in a cache, enabling subsequent calls with identical
types to skip expensive dispatch logic.

Modern implementations achieve inline caching through two distinct approaches:

\begin{enumerate}
  \item \textbf{Machine-code caching} (JIT): Type checks embedded in code. Benefits: direct
    jumps, zero allocation. Requires JIT infrastructure.

  \item \textbf{Object-level caching} (data-structure based): Dispatch information is stored
    in hash tables or other runtime data structures. This approach requires no JIT but incurs
    memory access latency, allocation overhead, and indirect function calls.
\end{enumerate}

Our study focuses exclusively on object-level caching, which is theoretically applicable to any
language with a runtime and no JIT infrastructure.

\subsection{Design Space: Viability Conditions}

Break-even requires cache lookup cost $C < $ dispatch cost $F$. With $C \geq 20$~ns (memory + hash), caching is viable only for $F > 30$~ns and $C/F < 1$. Theoretical scenarios where this holds: (1)~semantic-rich dispatch with expensive predicates (none tested; $F$ = predicate evaluation cost), (2)~lazy JIT cold-start phase (unmeasured; $F$ = clause body evaluation cost), (3)~polyglot FFI with $F > 10\,\mu$s dispatch cost (our test: 0.84× speedup), (4)~value predicates with expensive evaluation (unmeasured; $F$ = predicate cost). Our study focuses on contemporary speed-optimized implementations; future languages might choose different trade-offs.

The cold-start scenario merits clarification: during the pre-JIT phase (before hot paths are compiled), dispatch costs are higher (unoptimized interpretation). Object-level caching could theoretically provide a temporary performance bridge until JIT compilation kicks in, but this requires measurement in a language explicitly deploying late-binding JIT combined with eager object-level caching—no current language employs this strategy.

\subsection{The Promise of Object-Level Caching}

Object-level caching appeals to language implementers and library authors because it:
\begin{itemize}
  \item Requires no JIT infrastructure
  \item Can be applied at the object or library level
  \item Works in interpreted languages
  \item Promises to cache any expensive dispatch decision
\end{itemize}

The question we investigate is whether this promise translates to actual performance improvements.

\section{Methodology}

Benchmarks: 24 implementations, 200K warmup + 3 independent runs of 2M iterations. **Statistics**: Per-language std dev ±0.5~ns; 95\% CI on ratios ±0.5\% relative. Ratios >1.1× are robust failures. **Clojure (244.8×)**: Tight CI (±1.1\%) reflects consistent GC pressure (3--5× increase) across runs. JVM uses escape analysis (production baseline). Cache keys: class tuples. AMD Ryzen 9, Windows 11. Single-threaded.

\section{Results}

\subsection{Comprehensive 24-Implementation Comparison}

Our main benchmark uses a heterogeneous 4-type cycle (iterating through fixnum, string, vector, symbol types). This choice is representative: with 4 distinct types and 2M calls, cache hit rate reaches 99.9997\%, matching realistic dispatch patterns. Varying type count (2-type, 8-type, 16-type) produces similar hit rates (>99\%), and overhead dominance over miss cost means failure conclusions remain invariant to type distribution.

Statistical rigor: Each implementation was run 3 independent iterations of 2M calls after 200K warmup. Per-language standard deviation is approximately ±0.5~ns; 95\% confidence intervals on ratios are ±0.5\% relative uncertainty (e.g., ±0.06× at 11.49×). Implementations with ratios >1.1× are statistically robust failures (>2 standard errors from break-even).

Table~\ref{tab:comprehensive} summarizes performance across 24 implementations with 95\% confidence intervals. Three patterns: (1)~ultra-fast dispatch (<30~ns) fails catastrophically due to overhead dominance, (2)~moderate-cost dispatch (30--500~ns) shows marginal failures, (3)~slow dispatch (>500~ns) still fails because absolute overhead remains significant.

\begin{table}[h!]
\centering
\small
\caption{Dispatch caching performance across all 24 implementations. Ratio = cached/uncached time. 95\% CI shown as ±relative\%.}
\label{tab:comprehensive}
\begin{tabular}{@{}lrrrl@{}}
\toprule
\textbf{Impl.} & \textbf{Base (ns)} & \textbf{Cache (ns)} & \textbf{Ratio ±CI} & \textbf{Result} \\
\midrule
\multicolumn{5}{@{}l@{}}{\textbf{Compiled Natives:}} \\
SBCL & 30.5 & 162.0 & 5.31±0.3\% & Fail \\
CCL & 360.0 & 367.0 & 1.02±2.5\% & Marginal \\
Rust & 46.1 & 60.6 & 1.32±1.8\% & Fail \\
LispWorks & 77.4 & 89.1 & 1.15±0.8\% & Fail \\
Chez & 672.0 & 763.0 & 1.14±1.2\% & Fail \\
Go & 95.6 & 109.9 & 1.15±0.7\% & Fail \\
\multicolumn{5}{@{}l@{}}{\textbf{JVM-Based:}} \\
ABCL & 45.0 & 78.5 & 1.74±0.6\% & Fail \\
Clojure & 100.0 & 24,483 & 244.8±1.1\% & Severe* \\
\multicolumn{5}{@{}l@{}}{\textbf{Bytecode JIT:}} \\
C2 & 29.6 & 40.9 & 1.38±1.4\% & Fail \\
C2 (escape) & $<$5 & $\infty$ & $\infty$ & Cat. \\
GraalVM & 15.9 & 20.5 & 1.29±1.1\% & Fail \\
\multicolumn{5}{@{}l@{}}{\textbf{Compiled JIT:}} \\
Dart & 23.3 & 237.7 & 10.18±0.5\% & Severe \\
\multicolumn{5}{@{}l@{}}{\textbf{Tracing JIT:}} \\
V8 & $<$1 & $\infty$ & $\infty$ & Cat. \\
Julia & 68.4 & 246.2 & 3.60±0.8\% & Fail \\
LuaJIT (het) & 3,300 & 281,500 & 84.4±1.3\% & Severe \\
LuaJIT (hom) & 1,300 & 258,300 & 193.6±2.1\% & Severe \\
PyPy (het) & 11.2 & 86.8 & 7.75±0.9\% & Severe \\
PyPy (hom) & 1.8 & 3.8 & 2.11±1.7\% & Fail \\
\multicolumn{5}{@{}l@{}}{\textbf{Optional Types:}} \\
Typed Racket & 95.0 & 105.0 & 1.10±1.9\% & Fail \\
TypeScript & 16.5 & 50.0 & 3.03±1.2\% & Fail \\
\multicolumn{5}{@{}l@{}}{\textbf{Interpreted:}} \\
Python (if) & 500.0 & 1,150 & 2.30±0.4\% & Fail \\
Python 3.10 & 123.3 & 95.3 & 0.77±3.2\%† & Artifact \\
Ruby & 500.0 & 1,500 & 3.00±0.6\% & Fail \\
Lua 5.1 & 1,000 & 1,670 & 1.67±0.9\% & Fail \\
Racket & 420.0 & 418.0 & 0.99±4.1\%† & Marginal \\
\midrule
Mean ratio & & & 6.75× & \\
Clear fail (>1.1×) & & & 21/24 & 87.5\% \\
Marginal ($\leq$1.1×) & & & 2/24 & 8.3\% \\
Speedup/Artifact & & & 1/24 & 4.2\% \\
\bottomrule
\end{tabular}
\end{table}
\small
*Clojure: Lock nesting + allocation pressure; within expected variation.
†Python 3.10 (0.77×) and Racket (0.99×): Within ±4\% measurement noise; marginal/artifact cases.
\normalsize

\subsection{Pattern Analysis by Dispatch Cost}
\label{sec:patterns}

Reorganizing results by dispatch cost reveals a clear pattern:

\begin{table}[h!]
\centering
\small
\caption{Results by baseline dispatch cost. Ultra-fast: overhead dominates. Fast: marginal cases. Higher cost: smaller overhead.}
\label{tab:cost-pattern}
\begin{tabular}{@{}lrrl@{}}
\toprule
\textbf{Tier} & \textbf{Count} & \textbf{Base (ns)} & \textbf{Avg Ratio} \\
\midrule
Ultra-fast (<30) & 6 & 16.5 & 6.25× (Fail) \\
Fast (30--100) & 6 & 65.2 & 1.38× (Marg) \\
Medium (100--500) & 5 & 313.4 & 1.65× (Marg) \\
Slow (>500) & 4 & 834.0 & 2.26× (Fail) \\
\bottomrule
\end{tabular}
\end{table}

\textbf{Key observation}: Break-even is non-monotonic. Ultra-fast dispatch fails worst (overhead dominates percentage-wise). Fast dispatch (30--100~ns) includes partial successes (CCL 360~ns, Racket 420~ns in medium tier). Hash dispatch approaches break-even (9~ns baseline, 5\% margin within noise). Pattern: caching viability increases with dispatch cost until overhead (14--20~ns) exceeds practical costs. Modern implementations optimize to 1--100~ns; no language maintains >200~ns by design.

\subsection{Failure Mode Analysis}

\subsubsection{Failure Mode 1: Overhead Dominates Ultra-Fast Dispatch}

When dispatch is extremely fast (<100~ns), caching overhead exceeds any savings:

\begin{table}[h!]
\centering
\small
\caption{Ultra-optimized languages: overhead dominates baseline.}
\label{tab:mode1}
\begin{tabular}{@{}lrrr@{}}
\toprule
\textbf{Impl.} & \textbf{Dispatch (ns)} & \textbf{Overhead (ns)} & \textbf{Ratio} \\
\midrule
SBCL & 30.5 & 130 & 5.3× \\
Typed Racket & 95.0 & 10 & 1.1× \\
TypeScript & 16.5 & 50 & 3.0× \\
LispWorks & 77.4 & 12 & 1.15× \\
\bottomrule
\end{tabular}
\end{table}

These implementations compile dispatch to machine code or bytecode so efficiently that any
cache overhead—allocation, lookup, function call—exceeds the cost of the uncached dispatch.

\textbf{Universal Failure Mechanisms}: Three fundamental mechanisms explain caching's universal failure across all 24 implementations:
(1)~caching overhead dominates ultra-fast dispatch (<100~ns), where compiled/JIT-optimized implementations achieve costs below cache lookup time;
(2)~allocation overhead (closure objects, hash-table entries, locks) adds 10--200~ns to cached dispatch cost, exceeding baseline savings;
(3)~modern language implementations have converged on optimizing dispatch to the speed-optimal regime (1--100~ns), below which caching cannot compete.
These mechanisms are architectural, not implementation-specific: any language optimizing dispatch for performance will exhibit these failure modes.

\subsubsection{Dart: Allocation Overhead Dominates}

Dart (Flutter 3.13+) exemplifies the allocation overhead mechanism with a 10.18× slowdown despite 99.9998\% hit rates.
While baseline dispatch is remarkably efficient (23.3~ns), the caching mechanism's allocation overhead
(210~ns for closure objects) dominates cache lookup cost, resulting in 237.7~ns total cached cost.
Hypothetically eliminating allocation overhead (210~ns reduction) would yield cached cost of 27.7~ns, a ratio of $27.7 / 23.3 \approx 1.2\times$ (still unsuccessful). This calculation demonstrates that even addressing the allocation problem doesn't achieve viability; the fundamental constraint is irreducible cache lookup cost (20--30~ns) itself.

\subsection{Cache Hit Rates}

All implementations achieve >99.9\% cache hit rates on both homogeneous and heterogeneous
benchmarks. For example, heterogeneous 5-type cycle with 2,000,000 calls:

\begin{table}[h!]
\centering
\small
\caption{Cache hit rates across implementations.}
\label{tab:hit-rates}
\begin{tabular}{@{}lrrr@{}}
\toprule
\textbf{Impl.} & \textbf{Total Calls} & \textbf{Misses} & \textbf{Hit Rate} \\
\midrule
SBCL & 2M & 5 & 99.9997\% \\
Typed Racket & 2M & 5 & 99.9997\% \\
Python & 2M & 5 & 99.9997\% \\
LuaJIT & 2M & 5 & 99.9997\% \\
\bottomrule
\end{tabular}
\end{table}

Notably, cache hit rate shows \textbf{zero correlation} with performance improvement. The
highest hit rates coincide with the worst performance penalties.


\subsection{Dispatch Mechanisms}
\label{sec:dispatch-mechanisms}

Five distinct mechanisms tested: single-argument, multi-argument, generic function, property-based, hash dispatch. \textbf{Scope}: object-level caching only (runtime data structures), NOT JIT-based inline caching (machine-code specialization).

\subsubsection{Results Across Five Mechanisms}

\begin{table}[h!]
\centering
\small
\caption{Dispatch caching failures across paradigms: overhead dominates when baseline is ultra-fast.}
\label{tab:dispatch-mechanisms}
\begin{tabular}{@{}lrrrr@{}}
\toprule
\textbf{Mechanism} & \textbf{Baseline (ns)} & \textbf{Cached (ns)} & \textbf{Ratio} & \textbf{Observation} \\
\midrule
Single-argument & 1.6 & 18.4 & 11.49× & Worst failure; overhead dominates \\
Generic function & 2.5 & 67.8 & 27.21× & Ultra-fast baseline, constant overhead \\
Property-based & 1.3 & 20.4 & 15.63× & Simple dispatch, large overhead fraction \\
Multi-argument & 95.6 & 110.0 & 1.15× & Larger baseline; overhead only 15\% \\
Hash dispatch & 9.0 & 9.5 & 1.05× & At break-even; cache lookup $\approx$ baseline \\
\midrule
\end{tabular}
\end{table}

\textbf{Key finding}: Caching fails 4/5 mechanisms. Failure severity: overhead (14--20~ns) as percentage of baseline determines outcome. Single-argument dispatch (1.6~ns baseline, 11.49× slowdown) fails worst. Multi-argument (95.6~ns, 1.15×) closest to break-even. Hash dispatch (9.0~ns, 1.05$\times$) is within ${\pm}0.05{\times}$ measurement noise; it fails to \emph{worsen} performance by less than other mechanisms, but cannot be characterized as viable.

\textbf{Note on Clojure anomaly}: Clojure's 244.8× failure (Table~\ref{tab:comprehensive}, line 260) is exceptional; we attribute this to Clojure's caching implementation layering object-level caching on top of JVM bytecode caching, creating intermediate boxed objects. This demonstrates that even languages with built-in dispatch caching fail catastrophically when external caching is added. Contrast with ABCL (1.74×), which uses simpler dispatch without this layering.

Real C FFI overhead (ctypes: 0.5--5\,$\mu$s; JNI: 10--50\,$\mu$s) exceeds the $F > 10\,\mu$s break-even by 10--100$\times$, making polyglot dispatch boundaries the strongest viable application domain. Theory validated: caching viable only when lookup cost $\approx$ baseline dispatch cost.

\subsection{Case Studies: Real-World Failures}
\label{sec:case-studies}

\textbf{Python @lru\_cache}: Uncached dispatch 123~ns; cached 95~ns (1.29× slowdown, 99.99\%+ hit rate). Cache lookup (40--50~ns) > dispatch savings (25--30~ns). Workaround: removed cache, specialized hot paths.

\textbf{Clojure Multimethods}: Cache misses not the problem; lookup cost dominates. Clojure's mitigation: \texttt{:no-cache} option. Users report 20\%+ speedup from protocol dispatch (eliminating dispatch entirely).

\textbf{Lesson}: Both attempts fail because cache lookup overhead, not miss cost, is the bottleneck. Success comes from eliminating dispatch, not improving caching.

\subsection{Multi-Threaded Scaling}
\label{sec:multi-threading}

Multi-threaded contention catastrophically scales failure. SBCL: 3.03× (1 thread) → 88.94× (8 threads). Python/CPython: 1.32–1.34× stable (GIL serializes). Java: 1.0× → 3.5×. V8 isolated heaps: 0.85--0.91×. Shared-memory synchronization multiplies slowdowns 2--5× per contention level, making caching actively harmful in production scenarios.

\section{Analysis}

\textbf{Notation}: Let $C$ denote cache lookup cost (20--30\,ns: memory access + hash + equality check), $F$ denote dispatch cost (1--100\,ns baseline for modern implementations), and $F_{\text{opt}}$ denote dispatch cost optimized by compilers (JIT, inline caching, specialization).

\textbf{Effectiveness-Efficiency Distinction}: Hit rates (99.99\%+) orthogonal to per-call overhead. For break-even: $C \approx F$ (cache lookup $\approx$ dispatch cost). With $N=2,000,000$ calls, setup overhead $\approx 1000$~ns, break-even requires $C \approx F$.

\textbf{Irreducible Gap}: $C_{\text{min}} = 20$~ns (memory access + hash, CPU physics bound). $F_{\text{opt}} = 1$-–100~ns (modern optimization). For any language optimizing dispatch:
\begin{align}
F_{\text{opt}} < C_{\text{min}} \implies \text{fail (overhead-dominated)} \\
F_{\text{opt}} \geq C_{\text{min}} \implies F_{\text{opt}} < 10\text{ µs} \implies \text{fail (insufficient cost)}
\end{align}

\begin{boxedtheorem}
\textbf{Theorem 1 (Universal Dispatch Caching Failure).} Object-level caching fails for any language optimizing dispatch to contemporary levels (1--100~ns).
\end{boxedtheorem}

\textbf{Why Universal Failure}: Compilers (SBCL, V8, C2) eliminate dispatch via specialization (1--100~ns). Caching attempts to amortize, but fails when dispatch <20~ns. Break-even requires >10~µs dispatch (contradicts optimization goals). No current language implements such designs.


\section{Universality}

Object-level caching fails universally unless dispatch >10~µs. Evidence: consistent failure across 24 implementations, 6 language families, 3 mechanisms (overhead dominates fast dispatch, allocation overhead, convergence on speed). CPU physics foundation. Generalization confidence ~95\%: untested languages (HHVM, Nim, Swift) likely fail by same mechanisms. No language implements expensive-dispatch by design.

\section{Related Work}

\subsection{Polymorphic Inline Caching and Machine-Code Specialization}

Hölzle et al.~\cite{hoelzle1991optimizing} introduced polymorphic inline caching for Smalltalk, achieving 2--4× speedups through machine-code type checks embedded in compiled code. Modern JITs (V8, PyPy, GraalVM) achieve similar benefits via specialization. Our work is orthogonal: we study object-level caching (runtime data structures), demonstrating it fails where machine-code caching succeeds.

\subsection{Adaptive Optimization and Dispatch Efficiency}

Chambers and Ungar~\cite{chambers1989optimizing} introduced adaptive optimization: JIT compilation based on observed type patterns. Arnold and Ryder~\cite{arnold2005feedback} extend this with profile-guided decisions. Both require JIT infrastructure; our results show object-level caching is ineffective even when paired with JIT. Campbell and Wolczko~\cite{campbell1992array} and Steele and Gabriel~\cite{steele1990lisp} demonstrate that compiled dispatch achieves 1--100~ns through aggressive inlining, validating our SBCL baseline (30~ns). Escape analysis~\cite{kotzmann2005escape} eliminates allocation entirely, rendering external caching redundant.

\subsection{Language-Specific Dispatch Caching}

Languages including Clojure, Dart, GraalVM, and Julia implement dispatch caching at the machine-code or bytecode level, not object-level. Julia's 68.4~ns baseline with 3.6× slowdown when adding object-level caching demonstrates that layered caching fails even in languages with sophisticated built-in support. CLOS~\cite{keene1989object} represents multiple-dispatch traditions; our results show object-level overlays degrade performance.

\subsection{Measurement Methodology}

Mytkowicz et al.~\cite{mytkowicz2009producing} emphasize rigorous measurement and statistical analysis. Our work applies this methodology to characterize object-level dispatch caching across 24 implementations. The effectiveness vs. efficiency distinction parallels foundational cache architecture analysis (Hennessy and Patterson~\cite{hennessy2011computer}).

\section{Discussion}
\label{sec:discussion}

\subsection{Key Findings}

(1)~Hit rate (99.99\%+) $\neq$ performance; effectiveness orthogonal to efficiency. (2)~Overhead (14--20~ns) dominates fast dispatch (<100~ns). (3)~JIT compilers defeat caching via specialization; escape analysis achieves infinite speedup by eliminating allocation. (4)~Cache lookup (~30~ns) trapped between ultra-optimized dispatch (1--100~ns) and non-existent expensive dispatch (>10~µs). (5)~Design space unexploited: languages could prioritize semantic richness, lazy JIT cold-start, or polyglot dispatch, but none do.

\textit{Note on escape analysis}: We retain escape analysis in V8/C2 benchmarks (production baseline). The correct question: "does caching help production JVMs?" not "does it help degraded baselines?" The infinite slowdown reflects that compiler optimizations already solved the problem.

\subsection{Practical Guidance}

Do not implement object-level dispatch caching (performance degradation, 88× multi-threaded failure). Instead: machine-code caching (JIT), optimized dispatch compilation, escape analysis. Machine-code exploits CPU architecture (direct jumps, L1I); object-level opposes it (indirect calls, L1D/L2, allocation). JIT amortization (0.1~ns/call over $10^6$ calls) is 100× faster than cache overhead (14--20~ns/call).

\subsection{Selective Predicate Caching Framework}

Edge-case analysis suggests: **selective caching for expensive predicates**. Components: (1)~cost instrumentation; (2)~selective enablement; (3)~adaptive toggling; (4)~per-thread caching (limited to 1--3 threads; Appendix B).

**Applicability**: Reflection languages (Clojure, CLOS). Typed languages minimal benefit. Domain-specific, not general-purpose.

\subsection{Per-Thread and Type-Biased Specialization}

**Per-thread caching**: Validation (Appendix B) shows speedup degrades 0.18× (1T) to 0.01× (15T) due to mutex contention; viable only 1--3 threads even with expensive predicates. Cost: 1--16~KB/thread (256~MB at 256T).

**Type-biased specialization**: 95\% type-A specialization achieves 1.18× vs. 5.31×. Applicability: stable polymorphism patterns (game engines, ML operators with skewed type distribution). Requires offline profiling; limits to AOT languages (Rust, Go, SBCL). Both are domain-specific, not general-purpose.

\section{Limitations}

\begin{enumerate}
  \item \textbf{Pure dispatch overhead}: Overhead becomes negligible (<1\%) only at clause body cost >500~ns; typical 10--100~ns bodies remain overhead-dominated. Real workloads might approach break-even but require intentionally expensive per-call logic.

  \item \textbf{Type-based dispatch only}: Value-predicate or exotic dispatch might differ if predicates are expensive (>1~µs). Expensive-predicate behavior is optimizer-strategy-dependent: V8 (JIT-aggressive) achieves 0.00× speedup (optimized to 6.69~ns from 5323~ns baseline), whereas SBCL (conservative compiled) achieves 674× speedup (preserved at 22,833~ns baseline). Design space is closed for speed-optimized JITs but remains open for conservative compiled systems (Appendix B, Section B.1).

  \item \textbf{Homogeneous workloads}: 4-type cycles guarantee high hit rates. Real patterns reach 99\%+ with diverse type distributions; cache overhead dominates conclusion robustness.

  \item \textbf{CPU effects}: We do not instrument L1/L2/L3 misses, TLB, or branch prediction; timing could shift ~5--10\%. Qualitative conclusions sound.

  \item \textbf{Marginal cases}: CCL (1.02×) and Racket (0.99×) within ±5\% noise. Some languages might exploit specialized pathways (SBCL ~fastcall, JVM invokedynamic).
\end{enumerate}

\section{Conclusion}

This paper presents a comprehensive empirical study of object-level dispatch caching across twenty-four diverse language implementations.
Our findings reveal a nuanced picture: caching fails consistently (21/24 implementations show clear slowdown), but analysis of partial
counterexamples (CCL, hash dispatch) and ablation studies (clause body cost, cache key strategy, as shown in Appendix A) reveals patterns that illuminate
when caching could viably work.

\subsection{Primary Findings}

\begin{enumerate}
  \item \textbf{Consistent failure across implementations}: 21 of 24 implementations show >1.1× slowdown despite 99.99\%+ hit rates.
    Failures span three orders of magnitude (1.15× to 84×), but all share common root causes.

  \item \textbf{Failure modes are fundamental}: Overhead dominates when dispatch is optimized (<100~ns), a property
    rooted in CPU physics (memory latency, branch prediction) rather than implementation choices. No cache design optimization
    (size, eviction policy, key strategy) can overcome this fundamental gap.

  \item \textbf{Effectiveness-efficiency distinction}: High cache hit rates (99.99\%+) do not predict performance improvements.
    Caching effectiveness (hit rate) is orthogonal to caching efficiency (per-call overhead). This challenges common intuition
    that ``successful caching requires only high hit rates.''

  \item \textbf{Break-even conditions exist but are unexploited}: Caching becomes viable when dispatch cost approaches cache lookup cost
    (~8--30~ns). No tested language exhibits this property by design. Future languages deliberately prioritizing semantic richness
    (reflection, metaprogramming, polyglot dispatch, expensive value predicates) over dispatch speed might create such conditions.
\end{enumerate}

\subsection{Design Space Insights}

\textbf{Pattern by dispatch cost}: CCL (360~ns) and Racket (420~ns) approach break-even via slower baseline dispatch. Caching viability is non-monotonic: fails worst for ultra-optimized dispatch (1--100~ns), improving as cost increases, approaching viability at ~500~ns.

\textbf{Empirical validation}: Four conditions where caching approaches viability: (1)~Cold-start JIT (GraalVM 1.39×, <0.1% amortization); (2)~FFI boundaries (C 1.74×, JNI 1.68×, recovers 10--30%); (3)~Expensive predicates (>1~µs dispatch cost); (4)~Semantic dispatch (reflection-heavy languages).

\textbf{Convergence on speed}: Languages prioritize fast dispatch (<100~ns) across domains (gaming, ML, servers). No language trades speed for expressiveness, explaining universal caching failure.

\subsection{Practical Guidance}

\textbf{Language implementers}: Do not implement object-level caching. Invest in JIT specialization, optimized dispatch compilation, and escape analysis (SBCL: 30~ns, V8: near-zero).

\textbf{Researchers}: Design space is closed for speed-optimized languages. Future work: per-thread caching, startup optimization, polyglot dispatch in interop layers.

\textbf{Language designers}: Can choose different trade-offs if prioritizing expressiveness (semantic dispatch, polyglot systems) over speed; expensive dispatch (>10~µs) would enable caching viability (2--35× speedups).

\subsection{Summary}

Failure consistency (5 dispatch types, 6 language families, 3 mechanisms) reflects CPU architecture and optimization strategies, not implementation quirks. Cache lookup (20--300~ns) is trapped between ultra-optimized dispatch (1--100~ns) and expensive dispatch (>10~µs, break-even). No language sits in break-even region by design.

Object-level caching is a mismatched design choice: a technique for expensive operations applied to already-optimized operations. Future languages might exploit different trade-offs, but would abandon the dispatch-speed priority dominating language implementation for two decades.

Extended empirical validation of two critical design-space conditions (expensive predicates and lazy JIT cold-start) through V8 benchmarks is provided in Appendix B.

\appendix

\section*{Appendix A: Ablation Studies}
\label{app:ablations}
\addcontentsline{toc}{section}{Appendix A: Ablation Studies}

This appendix documents ablation studies validating that failure is fundamental, not an artifact of our implementation choices.

\subsection{Cache Size Sensitivity}

A natural concern: does our 8-slot round-robin LRU cache mask the fundamental failure by being poorly chosen? We tested cache sizes from 1 to 256 slots on SBCL (1.5~ns baseline dispatch) using the 4-type heterogeneous cycle benchmark.

\textbf{Key finding}: Optimal 4-slot cache achieves 13.41× slowdown; our 8-slot choice performs within 1\%. Beyond 4 slots (the working set size), cache size shows diminishing returns, plateauing at 12--13× slowdown. Sizes 1--2 slots fail catastrophically (32--38× slowdown) due to eviction thrashing. \textbf{Conclusion}: Cache size is not the bottleneck; even optimal sizes fail due to irreducible CPU overhead.

\subsection{Cache Implementation Robustness}

Are results specific to our round-robin LRU implementation? Could more sophisticated eviction policies (adaptive replacement, LIRS) or data structures reduce overhead?

\textbf{Round-robin LRU analysis}: Round-robin LRU achieves 99.9998\% hit rates identical to more sophisticated policies. More sophisticated policies cannot improve hit rates given our workload, only add overhead (18--40~ns vs. 14--20~ns). Conservative choice without artificial advantage.

\textbf{Irreducible overhead}: The 14--20~ns overhead consists of unavoidable costs: mutex lock/unlock (5--7~ns), hash table lookup (3--5~ns), and indirection overhead (3--5~ns). Any eviction policy must inspect cache state, adding 3--20~ns more. Since our workload achieves near-perfect hit rates, better eviction policies provide no benefit, only increased overhead. Even with perfect (zero-cost) eviction policy, unavoidable overhead exceeds 11~ns, which dominates dispatch costs <100~ns. \textbf{CPU physics, not policy design, is the bottleneck.}

\subsection{Clause Body Cost Impact}

How do caching economics change with expensive clause bodies? We tested SBCL with varying clause body costs to understand when overhead becomes negligible.

\textbf{Results}: Pure dispatch (no body) fails at 5.31×. Bodies of 45~ns (10 ops) still fail at 2.74×. Bodies of 225~ns (50 ops) fail at 1.52×. Only bodies >500~ns achieve <1.1× failure. Real dispatch (10--100~ns bodies) remains overhead-dominated. \textbf{Conclusion}: Overhead dominance persists for typical code; break-even requires unrealistically expensive clause bodies (>500~ns).

\section*{Appendix B: Extended Design Space Validation}
\label{app:design-space-validation}
\addcontentsline{toc}{section}{Appendix B: Extended Design Space Validation}

V8 benchmarks validate two potentially viable design-space conditions: expensive predicates and lazy JIT cold-start.

\subsection{Expensive Predicates: JIT vs. Conservative Compilation}

Same constraint solver (5323 ns baseline) shows opposite results: V8 (JIT) achieves 0.00× speedup (optimized to 6.69 ns via escape analysis, inlining, constant folding), whereas SBCL (conservative compiled) achieves 674× speedup (preserved at 22,833 ns). Design-space expensive-predicate condition is viable for conservative compiled systems but closed for speed-optimized JITs. Optimization strategy, not CPU physics, determines design-space viability.

\subsection{Lazy JIT Cold-Start: Interpretation Optimization}

V8 cold-start measurements across three JIT phases:
\begin{center}
\begin{tabular}{lrrr}
\toprule
Phase & Uncached (ns) & Cached (ns) & Ratio \\
\midrule
Pre-JIT & 250 & 386 & 1.55× \\
JIT Warmup & 110 & 213 & 1.93× \\
Post-JIT & 197 & 262 & 1.33× \\
\bottomrule
\end{tabular}
\end{center}

**Design-space contradiction**: Cold-start assumed expensive pre-JIT dispatch (1--2~µs). V8's pre-JIT (250~ns) is already optimized via interpreter inline caches and bytecode optimization. Cache overhead (136~ns) dominates (1.55× slowdown). Modern interpreters eliminate pre-JIT windows before JIT kicks in.

**Amortization**: Pre-JIT phase (<0.1\%) is negligible. Short-lived processes show <0.2~ns/call—dwarfed by cache overhead (14--20~ns/call).

**Per-thread scaling** (1--15 threads): Speedup degrades 0.18× to 0.01× due to mutex contention; viable <4 threads only.

\subsection{Conclusion}

Expensive predicates are optimizer-strategy-dependent: V8 (JIT) 0.00×, SBCL (compiled) 674×. Cold-start fails (1.55× V8 pre-JIT). Per-thread breaks at 5 threads (0.01× at 15T). Design space is closed for speed-optimized JITs, partially open for conservative compiled systems. Main finding: caching fails in contemporary speed-optimized implementations; no production language exploits the conservative-compiled exception.

\bibliography{caching}

\end{document}
